\documentclass[10pt,a4paper]{amsart}
\usepackage[margin=19mm]{geometry}
\usepackage{amsmath,amssymb,mathtools}
\usepackage[scr=rsfs,cal=euler]{mathalfa}
\usepackage{xfrac}
\usepackage[unicode,hidelinks,bookmarks=false]{hyperref}
\newcommand{\ie}{i.e.\ }
\newcommand{\cf}{cf.\ }
\newcommand{\resp}{respectively\ }
\newcommand{\Subsection}{\subsection}
\providecommand{\nobreakdash}{\nobreak\hbox{-}}
\setlength{\emergencystretch}{2em}
\multlinegap0pt
\usepackage{bm}
\usepackage{bbm}
\usepackage{dsfont}
\usepackage{xspace}
\usepackage{cancel}
\usepackage{mathtools}
\usepackage{amsthm}
\makeatletter
\@ifundefined{thm}{\newtheorem{thm}{Theorem}}{}
\makeatother
\newcommand{\vs}{\vspace{10pt}}
\title[Conjugacy Problem for Dehn Twists of Free Products of Free A]{Conjugacy Problem for Dehn Twists of Free Products of Free Abelian Groups}
\author{Amir Y. Weiss Behar, Chris Karpinski, Bratati Som}
\date{Source: arXiv:2606.02558v1 (CC BY 4.0)}
\begin{document}
\maketitle

\noindent\emph{Source and licence.} Conjugacy Problem for Dehn Twists of Free Products of Free Abelian Groups. Version arXiv:2606.02558v1 (CC BY 4.0), distributed under the Creative Commons Attribution 4.0 International licence, as recorded in the arXiv licence metadata for this exact version and shown on the abstract page https://arxiv.org/abs/2606.02558v1. Full rights notice: licenses/arxiv-2606.02558-CC-BY-4.0.txt.\par\smallskip

\noindent\textit{Editorial scope.} Counted unit: the Thm whose source label is \texttt{lemma: enlarging the splitting in separating case} in the article (source0.tex lines 272--275, source label \texttt{lemma: enlarging the splitting in separating case}), delivered together with the article's own complete original proof of it (source0.tex lines 277--355, one \texttt{proof} environment of 79 lines). No mathematics is written, repaired or re-derived: statement and proof are the article's own text line for line. Required context retained uncounted: none. Every definition, display and label of those lines is retained unchanged. Omitted from this package and not redistributed: 1--271, 276--276, 356--1314. Nothing inside a retained excerpt is omitted or altered: every retained body is delivered exactly as the article's lines are written, so no omission receipt of any kind accompanies this package. Citations: the retained excerpts carry no citation command at all, so no bibliography entry of the article is transcribed and no reference is redistributed. Reference adaptation: none was needed and none was made; every reference inside the delivered unit resolves inside it, so no unresolved reference or citation is produced. Printed slips kept literal: no printed word, symbol, hypothesis or number of the counted statement or of its proof was altered. Layout and class: the article's own class is \texttt{amsart} with packages including graphicx, amssymb, url, amsmath, cite, mathrsfs, hyperref, amsthm, tikz-cd, caption, float, inputenc, stmaryrd, tikz; the wrapper is the lane's pinned \texttt{amsart} editorial wrapper at 10pt on A4, which restates the theorem-like environments the retained text needs and adds the article's own macro definitions that the retained text needs (\texttt{\textbackslash vs}), with extra wrapper packages (bm, bbm, dsfont, xspace, cancel, mathtools). The delivered numbering is the unit's own. Assets: no retained span contains a figure environment, a graphics-inclusion command, a PDF-page inclusion, a table or a TikZ picture; no image file is copied into this package, the package declares no asset dependency and no image file is redistributed with it. Licence: version arXiv:2606.02558v1 of this article is distributed under the Creative Commons Attribution 4.0 International licence, as recorded in the arXiv licence metadata for this exact version and shown on the abstract page https://arxiv.org/abs/2606.02558v1. A full rights notice is delivered at licenses/arxiv-2606.02558-CC-BY-4.0.txt. Characters: every retained excerpt is pure ASCII, so the delivered engine, XeLaTeX with its default Latin Modern text font, reports no missing character.
\begin{thm}
    \label{lemma: enlarging the splitting in separating case}
    Let $G$ be a finitely generated group that splits as a free product of free abelian groups. Suppose that $G = A_0 *_{C_0} B_0$ where $C_0$ is non-trivial abelian. Fix some $c_{0} \in C_0 \setminus 1$ and let $\delta$ be a separating Dehn twist of $G$ associated with this splitting and $c_0 \in C_0$. Then there exists a splitting $G = A *_C B$ with $A = \mathrm{Fix}(\delta)$ (the fixed point set of $\delta$), $C \geq C_0, B \geq B_0$, $C \leq A, B$ and $C$ maximal abelian in $G$ (and hence in $A$ and $B$) and such that $\delta$ still acts as a separating Dehn twist of this splitting. 
\end{thm}

\begin{proof}
    First, note that $C_0$ is maximal abelian in $A_0$ or $B_0$. Indeed, if not, then taking $a_0 \in A_0 \setminus C_0$ and $b_0 \in B_0 \setminus C_0$ centralizing $C_0$, since $C_0 \neq 1$ and $G$ is commutative transitive (as it is CSA), we have that $a,b$ commute. However, this contradicts the normal form theorem, since there can be no relation between elements of $A_0 \setminus C_0$ and elements of $B_0 \setminus C_0$. 

    First, we will show that we can enlarge the splitting to a splitting $G = A *_C B$ where $C$ is malnormal in $A$ and $B$, and then we will show that in fact $A = \mathrm{Fix}(\delta)$. 

    Without loss of generality, suppose that $C_0$ is maximal abelian in $B_0$ and let $C_A$ be the centralizer of $C_0$ in $A_0$ (which is the maximal abelian subgroup of $A_0$ containing $C_0$). We can assume that $C_A \neq C_0$, since otherwise $C_0$ is already malnormal in both $A_0$ and $B_0$ and we are done. Let $A = A_0, B = \langle B_0, C_A \rangle$ and $C = C_A$. We claim that $G = A *_C B$, that $C$ is maximal abelian in both $A$ and $B$ and $\delta$ still acts as a separating Dehn twist of this splitting with respect to $c_0$. 

\vs

\underline{$C$ is maximal abelian in $A$}: Since $C$ is the centralizer of $C_0$ in $A$, we have that $C$ is maximal abelian in $A$.

\vs

\underline{$C$ is maximal abelian in $B$}: Suppose $g \in B$ is in the maximal abelian subgroup of $B$ containing $C$. Then $g$ commutes with some $c \in C \setminus C_0$. We can write $g = b_1 \cdots a_n$ for some $a_i \in C_A \setminus B_0 \subset A_0 \setminus B_0 = A_0 \setminus C_0$ and $b_i \in B_0 \setminus C_A \subset B_0 \setminus C_0$, with all letters non-trivial except possibly $b_1$ or $a_n$. We have $cgc^{-1}g^{-1} = 1$, which yields: 

$$ c (b_1 a_1 \cdots a_n) c^{-1} (b_1 a_1 \cdots b_n a_n)^{-1} = 1 \iff c b_1 a_1 \cdots b_n a_n c^{-1}a_n^{-1} b_n^{-1} \cdots a_1^{-1} b_1^{-1} = 1$$

Since $a_n \in C$, this becomes: 


$$ c b_1 a_1 \cdots b_n c^{-1} b_n^{-1} \cdots a_1^{-1} b_1^{-1} = 1$$

Since $c \in C \setminus C_0$, we have $c \in A_0 \setminus C_0$. If $n \geq 1$ and $g \notin C$, then $b_n \notin C_A$, so $b_n \in B_0 \setminus C_0$, We thus have an alternating word in $A_0 \setminus B_0$ and $B_0 \setminus C_0$ being the trivial element, contradicting the normal form theorem for $G = A_0 *_{C_0} B_0$. Therefore, $n = 1$ and $g = a_1 \in C_A$. We conclude that $C$ is maximal abelian, and hence malnormal, in $B$.

\vs 


\underline{$A \cap B = C$}: Note that $C \subseteq A,B$ and hence $C \subseteq A \cap B$. Let $g \in A \cap B$. Since $g \in B = \langle B_0, C_A \rangle$, we can write $g = a_1 b_1 \cdots a_n b_n$ for some $a_i \in C_A \setminus B_0 \subset A_0 \setminus B_0 = A_0 \setminus C_0$ and $b_i \in B_0 \setminus C_A \subset B_0 \setminus C_0$. Suppose for contradiction that $g \notin C$, so there is at least one letter $b_i$. We then have: 

$$1 = a_1 b_1 \cdots a_n b_n g^{-1}$$

Since $g \in A = A_0$, this is an alternating non-trivial word in $A_0 \setminus C_0$ and $B_0 \setminus C_0$, contradicting the normal form theorem for $G = A_0 *_{C_0} B_0$. Therefore, we must have $g \in C_A = C$. 

\vs

\underline{$G = A *_C B$}: Since $G = \langle A_0, B_0 \rangle$ and since $A \geq A_0, B \geq B_0$ we have $G = \langle A, B\rangle$. Since we showed above that $A \cap B = C$, it remains to show the normal form theorem for $A,B,C$. Suppose that: 

$$c a_1 b_1 \cdots a_n b_n = 1$$

for $c \in C, a_i \in A \setminus C, b_i \in B \setminus C$. 


Note that $A \setminus C \subset A_0 \setminus C_0$, so each $a_i \in A_0 \setminus C_0$. 

We can write each $b_i$ as $b_i = b_{i1}a_{i1} \cdots b_{im_i} a_{im_i}$ for $b_{ij} \in B_0 \setminus C_A$ and $a_{ij} \in C_A \setminus B_0$.  Substituting these expressions above, we obtain: 

$$ c a_1 (b_{11}a_{11} \cdots b_{1m_1} a_{1m_1})a_2 \cdots a_n (b_{n1}a_{n1} \cdots b_{nm_n} a_{nm_n}) = 1 $$

Note that $B_0 \setminus C_A \subset B_0 \setminus C_0$ and $C_A \setminus B_0 \subset A_0 \setminus C_0$, so each $b_{ij} \in B_0 \setminus C_0$ and $a_{ij} \in A_0 \setminus C_0$. Also, note that $ca_1 \notin C_0$, since otherwise $a_1 \in C$. Thus, we obtain a non-empty word in $A_0 \setminus C_0$ and $B_0 \setminus C_0$ equal to the identity in $G$, contradicting the normal form theorem for $G = A_0 *_{C_0} B_0$. We conclude that the word the $c a_1 b_1 \cdots a_n b_n$ must be empty, and so the normal form theorem holds for $A,B,C$ and thus that $G = A *_C B$. 

\vs 

\underline{$\delta$ still acts as a separating Dehn twist of $G = A *_C B$}: We still have $\delta$ being the identity on $A = A_0$. We need to show that $\delta$ acts by conjugation via $c_0$ on $B$. We know this is true on $B_0$. On $C_A$, $\delta$ fixes $C_A$ pointwise since $C_A \subset A$. Also, since $C_A$ is abelian and contains $c_0$, we have that $\delta$ acts as conjugation by $c_0$ on $C_A$ (which is trivial). Therefore, $\delta$ acts as conjugation by $c_0$ on $B$ as well. We conclude that $\delta$ acts as a separating Dehn twist of $G = A *_C B$. 

\vs 

Now we show that $A = \text{Fix}(\delta)$. Indeed, by definition $A \subseteq \text{Fix}(\delta)$. Let $h \in \text{Fix}(\delta)$. As $h \in A_0 *_{C_0} B_0$, it has a normal form, say, $h = c_h a_1b_1 \cdots a_nb_n$ where $c_h \in C_0$, each $a_i \in A_0 \setminus C_0$ and each $b_i \in B_0 \setminus C_0$. Now 
    \[
    c_h a_1(b_1)^{c_0} \cdots a_n(b_n)^{c_0} = \delta (h) = h = c_h a_1b_1 \cdots a_nb_n
    \implies (b_1)^{c_0} \cdots a_n(b_n)^{c_0} b_n^{-1}a_n^{-1} \cdots b_1^{-1} = 1
    \]

Then there must be cancellations, in particular $(b_n)^{c_0} b_n^{-1} \in C_0$ (otherwise it will contradict the fact that every reduced alternating word in $A_0 *_{C_0} B_0$ is non-trivial). This implies that $b_n c_0^{-1} b_n^{-1} \in C_0 \implies C_0^{b_n} \cap C_0\neq 1$. But $C_0$ is maximal abelian in $B_0$, hence malnormal in $B_0$. Thus, we obtain $b_n \in C_0$. Then $(b_n)^{c_0} b_n^{-1} = 1$ will reduce $(b_1)^{c_0} \cdots a_n(b_n)^{c_0} b_n^{-1}a_n^{-1} \cdots b_1^{-1} = 1$ to $(b_1)^{c_0} \cdots a_{n-1}(b_{n-1})^{c_0} b_{n-1}^{-1}a_{n-1}^{-1} \cdots b_1^{-1} = 1$. By the same argument as before, we can show that $b_{n-1} \in C_0$. Continuing the same way we get each $b_i$ is in $C_0$ implies hence $h \in A$. Hence $\text{Fix} (\delta) = A$.

\vs

We will next show that $C$ is maximal abelian in $G$. Let $C_G$ be the maximal abelian subgroup of $G$ containing $C$. We will show that $C_G \subseteq C$. Let $g \in C_G$. Then $g$ has a unique normal form, say, $g = c' a_1b_1 \cdots a_n b_n$, for $c' \in C_0$, $a_i \in A_0 \setminus C_0$ and $b_i \in B_0 \setminus C_0$. Since $g \in C_G$, we have that $g$ commutes with $c_0$, and hence $c_0 g c_0^{-1} = g$. This yields: 


$$ c' a_1^{c_0}b_1^{c_0} \cdots a_n^{c_0} b_n^{c_0} = c' a_1b_1 \cdots a_n b_n \implies a_1b_1 \cdots a_n b_n (b_n^{-1})^{c_0}(a_n^{-1})^{c_0} \cdots (b_1^{-1})^{c_0} (a_1^{-1})^{c_0} = 1$$


By using the same method for proving $\text{Fix} (\delta) = A$ above (which boiled down to centralizers of $C_0$ being maximal abelian and hence malnormal), we can deduce that each $a_i$ must be in the centralizer $C_A$ of $C_0$ in $A_0$ and each $b_i$ must be in the centralizer of $C_0$ in $B_0$, which is $C_0$ since $C_0$ is maximal abelian in $B_0$. This implies that $g \in A_0$. This proves the claim. 



Therefore, $\text{Fix}(\delta) = A_0 = A$ and $C = C_A = C_G$ is maximal abelian in $G$ and hence maximal abelian in both $A$ and $B$.

\end{proof}

\end{document}
